% GENERATED FILE. Edit canonical lessons, then run npm run generate:books.
% canonical-source-hash: fnv1a64:befdb6997d2c97b1
% canonical-lessons: IT-C337-condire, IT-C337-salare, IT-C337-insaporire, IT-C337-farcire, IT-C337-guarnire

\chapter{Il tocco finale}
\label{ch:it-il-tocco-finale}
\hlchaptermodality{337}

\begin{chapteropening}
\textbf{By the end of this chapter:} \emph{I can finish a dish: dressing, salting, flavouring, stuffing and garnishing.}

Complete the last lesson of chapter 337: I can finish a dish: dressing, salting, flavouring, stuffing and garnishing.
\end{chapteropening}

\section[condire]{condire — to dress, to season}

\label{lesson:IT-C337-condire}



\begin{quote}
Before the new one: say the Italian for the baking tray, then the Italian for the cooking.
\end{quote}

\subsection*{You'll want to know: condire}
\textbf{condire} — \textquotedblleft{}to dress, to season\textquotedblright{}.

\begin{quote}
Condisco l'insalata con olio e limone.
\end{quote}

\begin{quote}
\emph{I dress the salad with oil and lemon.}
\end{quote}

Present like \emph{capire}: \emph{condisco}.

\begin{cousinweb}
\emph{condire} is inherited from Latin \emph{condīre} \textquotedblleft{}to pickle, to season\textquotedblright{} (EtymDB).
\end{cousinweb}

\begin{grammarlens}[title={What you've built}]
One more word for telling, explaining and describing at B1.
\end{grammarlens}

\subsection*{Your turn}
\begin{itemize}
\raggedright
  \item \emph{Say it:} \emph{condire}
  \item \emph{Say it:} \emph{condire}, once more
  \item \emph{Say it:} a sentence of your own with \emph{condire}
  \item \emph{Recall:} say the Italian for the baking tray, then the Italian for the cooking, then say \emph{condire} again
\end{itemize}

\begin{tcolorbox}[breakable,colback=teal!4,colframe=teal!35!black,title={Before you move on}]
What does condire mean? (\textquotedblleft{}To dress, to season\textquotedblright{}.) Say it once more.
\end{tcolorbox}
\section[salare]{salare — to salt}

\label{lesson:IT-C337-salare}



\begin{quote}
Before the new one: say the Italian for the cooking, then the Italian for to dress.
\end{quote}

\subsection*{You'll want to know: salare}
\textbf{salare} — \textquotedblleft{}to salt\textquotedblright{}.

\begin{quote}
Hai salato l'acqua della pasta?
\end{quote}

\begin{quote}
\emph{Have you salted the pasta water?}
\end{quote}

\emph{Salato} — salty; also expensive: \emph{il conto è salato}.

\begin{cousinweb}
\emph{salare} is built on \emph{sale} \textquotedblleft{}salt\textquotedblright{}, from Latin \emph{sāl} (EtymDB).
\end{cousinweb}

\begin{grammarlens}[title={What you've built}]
One more word for telling, explaining and describing at B1.
\end{grammarlens}

\subsection*{Your turn}
\begin{itemize}
\raggedright
  \item \emph{Say it:} \emph{salare}
  \item \emph{Say it:} \emph{salare}, once more
  \item \emph{Say it:} a sentence of your own with \emph{salare}
  \item \emph{Recall:} say the Italian for the cooking, then the Italian for to dress, then say \emph{salare} again
\end{itemize}

\begin{tcolorbox}[breakable,colback=teal!4,colframe=teal!35!black,title={Before you move on}]
What does salare mean? (\textquotedblleft{}To salt\textquotedblright{}.) Say it once more.
\end{tcolorbox}
\section[insaporire]{insaporire — to flavour, to give taste to}

\label{lesson:IT-C337-insaporire}



\begin{quote}
Before the new one: say the Italian for to dress, then the Italian for to salt.
\end{quote}

\subsection*{You'll want to know: insaporire}
\textbf{insaporire} — \textquotedblleft{}to flavour, to give taste to\textquotedblright{}.

\begin{quote}
Le erbe insaporiscono la carne.
\end{quote}

\begin{quote}
\emph{The herbs flavour the meat.}
\end{quote}

Built on \emph{sapore}, ‘flavour’.

\begin{cousinweb}
\emph{insaporire} is built on \emph{sapore} \textquotedblleft{}flavour\textquotedblright{}, from Latin \emph{sapor} \textquotedblleft{}taste\textquotedblright{} (Lewis \& Short).
\end{cousinweb}

\begin{grammarlens}[title={What you've built}]
One more word for telling, explaining and describing at B1.
\end{grammarlens}

\subsection*{Your turn}
\begin{itemize}
\raggedright
  \item \emph{Say it:} \emph{insaporire}
  \item \emph{Say it:} \emph{insaporire}, once more
  \item \emph{Say it:} a sentence of your own with \emph{insaporire}
  \item \emph{Recall:} say the Italian for to dress, then the Italian for to salt, then say \emph{insaporire} again
\end{itemize}

\begin{tcolorbox}[breakable,colback=teal!4,colframe=teal!35!black,title={Before you move on}]
What does insaporire mean? (\textquotedblleft{}To flavour, to give taste to\textquotedblright{}.) Say it once more.
\end{tcolorbox}
\section[farcire]{farcire — to stuff, to fill}

\label{lesson:IT-C337-farcire}



\begin{quote}
Before the new one: say the Italian for to salt, then the Italian for to flavour.
\end{quote}

\subsection*{You'll want to know: farcire}
\textbf{farcire} — \textquotedblleft{}to stuff, to fill\textquotedblright{}.

\begin{quote}
Ho farcito i peperoni con riso e carne.
\end{quote}

\begin{quote}
\emph{I stuffed the peppers with rice and meat.}
\end{quote}

Cakes too: \emph{farcire una torta con la crema}.

\begin{cousinweb}
\emph{farcire} is inherited from Latin \emph{farcīre} \textquotedblleft{}to cram, to stuff\textquotedblright{} (EtymDB).
\end{cousinweb}

\begin{grammarlens}[title={What you've built}]
One more word for telling, explaining and describing at B1.
\end{grammarlens}

\subsection*{Your turn}
\begin{itemize}
\raggedright
  \item \emph{Say it:} \emph{farcire}
  \item \emph{Say it:} \emph{farcire}, once more
  \item \emph{Say it:} a sentence of your own with \emph{farcire}
  \item \emph{Recall:} say the Italian for to salt, then the Italian for to flavour, then say \emph{farcire} again
\end{itemize}

\begin{tcolorbox}[breakable,colback=teal!4,colframe=teal!35!black,title={Before you move on}]
What does farcire mean? (\textquotedblleft{}To stuff, to fill\textquotedblright{}.) Say it once more.
\end{tcolorbox}
\section[guarnire]{guarnire — to garnish}

\label{lesson:IT-C337-guarnire}



\begin{quote}
Before the new one: say the Italian for to flavour, then the Italian for to stuff.
\end{quote}

\subsection*{You'll want to know: guarnire}
\textbf{guarnire} — \textquotedblleft{}to garnish\textquotedblright{}.

\begin{quote}
Guarnisci il piatto con una fetta di limone.
\end{quote}

\begin{quote}
\emph{Garnish the dish with a slice of lemon.}
\end{quote}

Present like \emph{capire}: \emph{guarnisco}.

\begin{cousinweb}
\emph{guarnire} is usually traced to a Germanic verb \emph{warnjan}; EtymDB lists it beside Old French \emph{guarnir} \textquotedblleft{}to equip, to protect\textquotedblright{} and French \emph{garnir} \textquotedblleft{}to furnish, to garnish\textquotedblright{}.
\end{cousinweb}

\begin{grammarlens}[title={What you've built}]
One more word for telling, explaining and describing at B1.
\end{grammarlens}

\subsection*{Your turn}
\begin{itemize}
\raggedright
  \item \emph{Say it:} \emph{guarnire}
  \item \emph{Say it:} \emph{guarnire}, once more
  \item \emph{Say it:} a sentence of your own with \emph{guarnire}
  \item \emph{Recall:} say the Italian for to flavour, then the Italian for to stuff, then say \emph{guarnire} again
\end{itemize}

\begin{tcolorbox}[breakable,colback=teal!4,colframe=teal!35!black,title={Before you move on}]
What does guarnire mean? (\textquotedblleft{}To garnish\textquotedblright{}.) Say it once more.
\end{tcolorbox}
